\documentclass[10pt]{article}
\usepackage[verbose, a4paper, hmargin=2.5cm, vmargin=2.5cm]{geometry}

\usepackage{fontspec}
\usepackage{ctex}
\usepackage{paratype}

\usepackage{amsmath}
\usepackage{amssymb}
\usepackage{amsfonts}
\usepackage{esint}
\usepackage{stmaryrd}

\usepackage[mathbf=sym]{unicode-math}
\setmainfont{STIX Two Text}
\setmathfont{STIX Two Math}

\setsansfont{Arial}
\setmonofont{Courier New}

\usepackage{makecell}
\usepackage{wasysym}
\usepackage{pifont}
\usepackage[export]{adjustbox}
\usepackage{graphicx}
\usepackage{mdframed}
\usepackage{booktabs,array,multirow}
\usepackage{tabularx}
\usepackage{hyperref}
\hypersetup{colorlinks=true, linkcolor=blue, filecolor=magenta, urlcolor=cyan,}
\urlstyle{same}
\usepackage[most]{tcolorbox}
\definecolor{mygray}{RGB}{240,240,240}
\tcbset{
  colback=mygray,
  boxrule=0pt,
}
\graphicspath{ {./images/} }
\newcommand{\HRule}{\begin{center}\rule{0.9\linewidth}{0.2mm}\end{center}}
\newcommand{\customfootnote}[1]{
  \let\thefootnote\relax\footnotetext{#1}
}
\newcommand{\lt}{\mathrel{<}}
\newcommand{\gt}{\mathrel{>}}
\newcommand*\circled[1]{\tikz[baseline=(char.base)]{\node[shape=circle,draw,inner sep=1.5pt] (char) {\small #1};}}
\setcounter{MaxMatrixCols}{256}
\begin{document}
Lemma 1.7.1 If \(X\) is regular with valency \(k\) , then \(L\left( X\right)\) is regular with valency \({2k} - 2\) .

\begin{tcolorbox}\relax
引理 1.7.1 若 \(X\) 为度数为 \(k\) 的正则图，则 \(L\left( X\right)\) 为度数为 \({2k} - 2\) 的正则图。
\end{tcolorbox}

Each vertex in \(X\) determines a clique in \(L\left( X\right)\) : If \(x\) is a vertex in \(X\) with valency \(k\) , then the \(k\) edges containing \(x\) form a \(k\) -clique in \(L\left( X\right)\) . Thus if \(X\) has \(n\) vertices, there is a set of \(n\) cliques in \(L\left( X\right)\) with each vertex of \(L\left( X\right)\) contained in at most two of these cliques. Each edge of \(L\left( X\right)\) lies in exactly one of these cliques. The following result provides a useful converse:

\begin{tcolorbox}\relax
\(X\) 中的每个顶点在 \(L\left( X\right)\) 中确定一个团:若 \(x\) 是 \(X\) 中度数为 \(k\) 的顶点，则包含 \(x\) 的那 \(k\) 条边在 \(L\left( X\right)\) 中构成一个 \(k\) -团。因此若 \(X\) 有 \(n\) 个顶点，则在 \(L\left( X\right)\) 中存在一组 \(n\) 个团，且 \(L\left( X\right)\) 的每个顶点至多属于其中两个团。 \(L\left( X\right)\) 的每条边恰好属于这些团中的一个。下述结果给出了一个有用的逆定理:
\end{tcolorbox}

Theorem 1.7.2 A nonempty graph is a line graph if and only if its edge set can be partitioned into a set of cliques with the property that any vertex lies in at most two cliques.

\begin{tcolorbox}\relax
定理 1.7.2 非空图当且仅当其边集可划分为若干团，且任一顶点至多属于两团，这时该图是线图。
\end{tcolorbox}

If \(X\) has no triangles (that is, cliques of size three), then any vertex of \(L\left( X\right)\) with at least two neighbours in one of these cliques must be contained in that clique. Hence the cliques determined by the vertices of \(X\) are all maximal.

\begin{tcolorbox}\relax
若 \(X\) 无三角(即无三元团)，则 \(L\left( X\right)\) 中在某一团中至少有两邻点的任何顶点必然属于该团。因此由 \(X\) 的顶点确定的团均为极大团。
\end{tcolorbox}

It is both obvious and easy to prove that if \(X \cong  Y\) , then \(L\left( X\right)  \cong  L\left( Y\right)\) . However, the converse is false: \({K}_{3}\) and \({K}_{1,3}\) have the same line graph, namely \({K}_{3}\) . Whitney proved that this is the only pair of connected counterexamples. We content ourselves with proving the following weaker result.

\begin{tcolorbox}\relax
显然且易证若 \(X \cong  Y\) ，则 \(L\left( X\right)  \cong  L\left( Y\right)\) 。然而逆命题不成立: \({K}_{3}\) 与 \({K}_{1,3}\) 有相同的线图，即 \({K}_{3}\) 。惠特尼证明这是唯一的一对连通反例。我们满足于证明以下较弱的结果。
\end{tcolorbox}

Lemma 1.7.3 Suppose that \(X\) and \(Y\) are graphs with minimum valency four. Then \(X \cong  Y\) if and only if \(L\left( X\right)  \cong  L\left( Y\right)\) .

\begin{tcolorbox}\relax
引理 1.7.3 设 \(X\) 与 \(Y\) 为最小度为四的图。则 \(X \cong  Y\) 当且仅当 \(L\left( X\right)  \cong  L\left( Y\right)\) 。
\end{tcolorbox}

Proof. Let \(C\) be a clique in \(L\left( X\right)\) containing exactly \(c\) vertices. If \(c > 3\) , then the vertices of \(C\) correspond to a set of \(c\) edges in \(X\) , meeting at a common vertex. Consequently, there is a bijection between the vertices of \(X\) and the maximal cliques of \(L\left( X\right)\) that takes adjacent vertices to pairs of cliques with a vertex in common. The remaining details are left as an exercise.

\begin{tcolorbox}\relax
证明。令 \(C\) 为 \(L\left( X\right)\) 中恰含 \(c\) 个顶点的一个团。若 \(c > 3\) ，则 \(C\) 的顶点对应于 \(X\) 中在同一顶点相交的 \(c\) 条边。因此存在一个在相邻顶点对应有公共顶点的团对之间建立一一对应的映射，将 \(X\) 的顶点与 \(L\left( X\right)\) 的极大团一一对应。其余细节留作练习。
\end{tcolorbox}

There is another interesting characterization of line graphs:

\begin{tcolorbox}\relax
还有另一种关于线图的有趣刻画:
\end{tcolorbox}

Theorem 1.7.4 A graph \(X\) is a line graph if and only if each induced subgraph of \(X\) on at most six vertices is a line graph.

\begin{tcolorbox}\relax
定理 1.7.4 图 \(X\) 当且仅当其任一最多含六个顶点的诱导子图都是线图。
\end{tcolorbox}

Consider the set of graphs \(X\) such that

\begin{tcolorbox}\relax
考虑满足下列条件的图集 \(X\) :
\end{tcolorbox}
\section*{Exercises}

\begin{tcolorbox}\relax
\section*{练习}
\end{tcolorbox}

1. Let \(X\) be a graph with \(n\) vertices. Show that \(X\) is complete or empty if and only if every transposition of \(\{ 1,\ldots ,n\}\) belongs to \(\operatorname{Aut}\left( X\right)\) .

\begin{tcolorbox}\relax
1. 设 \(X\) 为有 \(n\) 个顶点的图。证明当且仅当每个 \(\{ 1,\ldots ,n\}\) 的换位都属于 \(\operatorname{Aut}\left( X\right)\) 时， \(X\) 是完全图或空图。
\end{tcolorbox}

2. Show that \(X\) and \(\bar{X}\) have the same automorphism group, for any graph \(X\) .

\begin{tcolorbox}\relax
2. 证明对任意图 \(X\) ， \(X\) 与 \(\bar{X}\) 具有相同的自同构群。
\end{tcolorbox}

3. Show that if \(x\) and \(y\) are vertices in the graph \(X\) and \(g \in  \operatorname{Aut}\left( X\right)\) , then the distance between \(x\) and \(y\) in \(X\) is equal to the distance between \({x}^{g}\) and \({y}^{g}\) in \(X\) .

\begin{tcolorbox}\relax
3. 证明若 \(x\) 与 \(y\) 是图 \(X\) 中的顶点且 \(g \in  \operatorname{Aut}\left( X\right)\) ，则在 \(X\) 中 \(x\) 与 \(y\) 的距离等于在 \(X\) 中 \({x}^{g}\) 与 \({y}^{g}\) 的距离。
\end{tcolorbox}

4. Show that if \(f\) is a homomorphism from the graph \(X\) to the graph \(Y\) and \({x}_{1}\) and \({x}_{2}\) are vertices in \(X\) , then

\begin{tcolorbox}\relax
4. 证明若 \(f\) 是从图 \(X\) 到图 \(Y\) 的同态，且 \({x}_{1}\) 与 \({x}_{2}\) 是 \(X\) 中的顶点，则
\end{tcolorbox}

\[
{d}_{X}\left( {{x}_{1},{x}_{2}}\right)  \geq  {d}_{Y}\left( {f\left( {x}_{1}\right) ,f\left( {x}_{2}\right) }\right) .
\]

5. Show that if \(Y\) is a subgraph of \(X\) and \(f\) is a homomorphism from \(X\) to \(Y\) such that \(f \upharpoonright  Y\) is a bijection, then \(Y\) is a retract.

\begin{tcolorbox}\relax
5. 证明若 \(Y\) 是 \(X\) 的一个子图，且存在从 \(X\) 到 \(Y\) 的同态 \(f\) 使得 \(f \upharpoonright  Y\) 为双射，则 \(Y\) 是一个收缩子图(retract)。
\end{tcolorbox}

6. Show that a retract \(Y\) of \(X\) is an induced subgraph of \(X\) . Then show that it is isometric, that is, if \(x\) and \(y\) are vertices of \(Y\) , then \({d}_{X}\left( {x,y}\right)  = {d}_{Y}\left( {x,y}\right) .\)

\begin{tcolorbox}\relax
6. 证明 \(X\) 的一个收缩子图 \(Y\) 是 \(X\) 的诱导子图。然后证明它是等距的，即若 \(x\) 与 \(y\) 是 \(Y\) 的顶点，则 \({d}_{X}\left( {x,y}\right)  = {d}_{Y}\left( {x,y}\right) .\)
\end{tcolorbox}

7. Show that any edge in a bipartite graph \(X\) is a retract of \(X\) .

\end{document}
